\chapter{Python Application Development}
\label{chap:python}

\section{Project Structure}

\begin{lstlisting}[language=bash, caption={Project Directory Structure}]
panic-selling-detection/
├── src/
│   ├── config.py                          # DB connection, paths, constants
│   ├── data_loader.py                     # Fetch market data via vnstock
│   ├── indicator.py                       # RSI, MA, feature engineering
│   ├── investor_data_generation.py        # Generate 1000 investor profiles
│   ├── seed_market_data.py                # Push market data to MySQL
│   ├── compute_behavior_signals.py        # Compute PanicScore → BehaviorSignals
│   ├── transaction_data_generation/
│   │   ├── investor.py                    # InvestorAgent behavioral logic
│   │   ├── portfolio.py                   # Portfolio state management
│   │   ├── prepare_market.py              # Market regime classification
│   │   └── simulation.py                  # Main simulation engine
│   └── webapp/
│       ├── app.py                         # Flask application
│       └── templates/
│           ├── index.html                 # Dashboard
│           ├── investor.html              # Investor profile page
│           └── warnings.html             # Warning log page
├── data/
│   ├── raw/                               # CSV outputs from simulation
│   └── processed/
├── database/
│   └── schema.sql                         # Full DDL + advanced objects
└── notebooks/
    └── test_data.ipynb                    # EDA notebook
\end{lstlisting}

\section{Database Connection}

\begin{lstlisting}[style=pythonstyle, caption={Database Connection via SQLAlchemy (config.py)}]
from sqlalchemy import create_engine

DB_CONNECTION_STR = (
    f"mysql+pymysql://{DB_USER}:{DB_PASSWORD}"
    f"@{DB_HOST}:{DB_PORT}/{DB_NAME}?charset=utf8mb4"
)
engine = create_engine(DB_CONNECTION_STR)
\end{lstlisting}

\section{Behavioral Analytics Pipeline}

The pipeline processes raw transaction data into Panic Scores through five stages.

% TODO: Add pipeline diagram
% \begin{figure}[H]
%     \centering
%     \includegraphics[width=0.95\textwidth]{figures/pipeline.png}
%     \caption{Behavioral Analytics Pipeline}
% \end{figure}

\subsection{Feature Engineering}

Three behavioral features are computed per investor per observation window:

\begin{enumerate}
    \item \textbf{DrawdownLevel}: Maximum percentage loss from portfolio peak within the past 30 days, computed via \texttt{fn\_rolling\_drawdown}.
    \item \textbf{SellSpike}: Ratio of SELL transactions executed during PANIC regime to total transactions, computed via \texttt{fn\_sell\_intensity}.
    \item \textbf{LossSensitivity}: Speed and magnitude of sell response to negative returns, measured by proportion of PANIC-labeled sell transactions.
\end{enumerate}

\subsection{Panic Score Computation}

\begin{lstlisting}[style=pythonstyle, caption={Panic Score computation script (compute\_behavior\_signals.py)}]
def compute_behavior_signals(engine, obs_date):
    """
    Compute BehaviorSignals for all investors for a given observation date.
    Calls MySQL stored procedure which internally uses UDFs and fires trigger.
    """
    with engine.begin() as conn:
        investors = pd.read_sql(
            "SELECT Investor_ID FROM investors", conn
        )
        for inv_id in investors['Investor_ID']:
            conn.execute(
                text("CALL sp_refresh_behavior_signals(:inv, :dt)"),
                {"inv": inv_id, "dt": obs_date}
            )
    print(f"Computed signals for {len(investors)} investors on {obs_date}")
\end{lstlisting}

\section{Flask Web Application}

The web application exposes three main routes.

\subsection{Dashboard (/)}

Displays aggregate statistics: distribution of Panic Scores across all investors, count of High/Medium/Low warnings, and top 10 highest-risk investors.

\subsection{Investor Profile (/investor/\textless id\textgreater)}

Accepts an investor ID and returns the full output format specified by the project requirements:

\begin{lstlisting}[style=pythonstyle, caption={Investor profile route}]
@app.route('/investor/<investor_id>')
def investor_profile(investor_id):
    query = """
        SELECT w.InvestorID, w.PanicLevel, w.Confidence, w.KeySignals,
               bs.PanicScore, bs.DrawdownLevel, bs.SellSpike
        FROM vw_warning_dashboard w
        JOIN BehaviorSignals bs ON w.InvestorID = bs.InvestorID
        WHERE w.InvestorID = :inv_id
        ORDER BY w.WarningDate DESC LIMIT 1
    """
    result = pd.read_sql(text(query), engine, params={"inv_id": investor_id})
    return render_template('investor.html', data=result.to_dict('records'))
\end{lstlisting}

\subsection{Warning Log (/warnings)}

Queries \texttt{vw\_warning\_dashboard} and displays all warnings filterable by panic level.

% TODO: Add screenshots of each page
% \begin{figure}[H]
%     \centering
%     \includegraphics[width=0.9\textwidth]{figures/dashboard_screenshot.png}
%     \caption{Web Application Dashboard}
% \end{figure}
